% sup4_06_S21a: Appendix on Ba3 Prop. 1, part 4, step S21, I: the identities of the renormalisation step (degree formula and the
% divergent types, C-parity, Taylor's formula, the max-scale regrouping, rescaling, the Ward-Takahashi identities at
% full lines, lattice reflections). Input from a08_balaban_prop1.tex.

\subsubsection{S21, I: the identities of the renormalisation step, \cite[pp.~434--445]{Ba3}}\label{sup4:ssec:S21a}

\cite[p.~434]{Ba3} states \cite[(3.5)]{Ba3}: for each class $G_{\mathrm{ren}}$ (after S13--S20, so with scalar lines $G_k(0)$ on $\eta\Z^d$, photon lines on the torus, vertices localised in $\square_1$ and all other objects external),
\[
\sum_{G\in G_{\mathrm{ren}}}\sum_{\tilde l}\sum_{j\in J(\tilde l)}E(G(j))=\sum_{G'\in G'_{\mathrm{fin}}}\sum_{\tilde l'}\sum_{j\in J(\tilde l')}\gamma(G',\tilde l')E'(G'(j)),
\]
with generalised graphs $G'$ and sets of orderings (indicator $\gamma$) along which every subgraph of the sequence has positive degree; \cite[Prop.~2.2]{Ba3} (\cref{sup2:thm:S6}) then gives \cite[(3.3)]{Ba3}. Ba{\l}aban writes out the lowest-order members of each divergent type and says ``the same procedure'' \cite[p.~438]{Ba3}, ``the same applies'' \cite[p.~437]{Ba3}, ``in the way described several times'' \cite[p.~444]{Ba3} for the rest. This subsection and the next prove the identities and the one-loop classes which the renormalisation theorem of \cref{app:a06_BALtorus} (Step~3(D$_4$), \cref{thm:a06d:S49}) uses; the two-loop remainder terms, the transport of legs across the subgraphs \cite[(2.4)]{Ba3} and the induction of \cite[pp.~444--445]{Ba3} are proved there. Throughout $d=3$, and $\alpha\in(0,\frac12)$ is Ba{\l}aban's Taylor exponent \cite[p.~435]{Ba3}, so that a loss of $\alpha$ never makes a half-integer degree non-positive. Every Ward--Takahashi identity below is applied to full lines (cumulative propagators summed over all internal scales below the external scale, then replaced by the free $C^\xi$), never to a fixed ordering or a fixed scale assignment, and every remainder of the replacement is bounded explicitly.

\begin{lemma}[degree formula and the divergent types]\label{sup4:lem:S211}
Let $G$ be connected and built from \cite[(1.6), (1.8), (1.10)]{Ba3} only. Let $m_v$ be the number of vector legs at $v$ (internal and external), $\delta_v\in\{0,1\}$ with $\delta_v=1$ iff the differentiated scalar leg of a \cite[(1.8)]{Ba3}-vertex is internal, $V_6$ the number of \cite[(1.6)]{Ba3}-vertices, and $E=E_s+E_v$ the number of external scalar and vector legs. Then
\[
D(G)=V_6+\sum_{(1.8)}\big(1+\tfrac{m_v}2-\delta_v\big)+\sum_{(1.10)}\tfrac{m_v}2+\tfrac E2-3 .
\]
The vertex weights are $\ge\frac12$ for \cite[(1.8)]{Ba3} (equality iff $m_v=1$, $\delta_v=1$) and $\ge1$ for \cite[(1.6), (1.10)]{Ba3}; $E_s$ is even. Mass-counterterm vertices count as the graph they replace \cite[(2.3)]{Ba3}, and graphs containing a block vertex \cite[(1.13)--(1.15)]{Ba3}, an $R$-vertex or a finite counterterm have $D>0$ (\cref{sup2:cor:S7class}). Hence $D(G)\le0$ only in the following cases: $E\ge5$ never; $E=4$ exactly the graph \cite[(2.4)]{Ba3}; $E=3$ with $(E_s,E_v)=(0,3)$, which vanishes by \cref{sup4:lem:S212}, or $(2,1)$, the $\phi\phi A$ class \cite[(2.18)--(2.21)]{Ba3}; $E=2$ with $(2,0)$, the scalar self-energies, $D=-1$ only for \cite[(3.6)]{Ba3} (the two tadpoles and \cite[(3.8)]{Ba3}), and $(0,2)$, the vector self-energies, $D=-1$ only for the four graphs \cite[(3.25)]{Ba3}; $E=1$ with $(0,1)$ vanishes by \cref{sup4:lem:S212}, and $(1,0)$ is impossible; $E=0$ the vacuum graphs. Divergent graphs with $E\ge2$, in particular every proper divergent subgraph, have at most four vertices; vacuum graphs can have up to six (six \cite[(1.8)]{Ba3}-vertices with $m=1$ on one scalar loop with three photon lines have $D=0$ at order $e^6$), which matches the vacuum counterterms up to $\alpha+2\beta\le6$ of \cite[p.~418]{Ba3}.
\end{lemma}

\begin{proof}
$D_G(v)=(\eta\text{-power})-(\text{internal legs})/2-\delta_v$ in $d=3$, with $\eta$-powers $3$, $3+m-1$, $3+m-2$ for \cite[(1.6), (1.8), (1.10)]{Ba3} (\cref{sup2:lem:S7vert}). Use $\sum_v(\text{internal legs})=\sum_v(\text{legs})-E$ with $4$, $2+m$, $2+m$ legs; then $D=\sum_vD_G(v)-3$ gives the formula. $D\le0$ forces $\sum_vw_v\le3-E/2$, and $w\ge\frac12$ bounds the number of vertices by $6-E$: by $4$ for $E\ge2$ and by $6$ for $E=0$. $E=4$, $\sum w\le1$: a single vertex with a self-line has $E\le2$; two vertices need $w=\frac12$ each, i.e.\ \cite[(1.8)]{Ba3} with $m=1$, $\delta=1$, joined by the scalar line that carries both derivatives, which is \cite[(2.4)]{Ba3}. $E=5$, $\sum w\le\frac12$: one \cite[(1.8)]{Ba3}-vertex with $\delta=1$ needs an internal scalar line, hence a second vertex or a self-line, both impossible. $E=2$, $D=-1$, $\sum w=1$: a single \cite[(1.6)]{Ba3} with a self-line; a single \cite[(1.10)]{Ba3} with $m=2$ (the photon tadpole); or two \cite[(1.8)]{Ba3} with $m=1$, $\delta=1$, i.e.\ \cite[(3.8)]{Ba3} for $E_s=2$ and the bubbles of \cite[(3.25)]{Ba3} for $E_v=2$; a single \cite[(1.8)]{Ba3} with $m=2$ has $\delta=0$ when both scalar legs are external ($w=2$) and $\delta=1$ when they form a self-line ($E_v=2$, the fourth graph of \cite[(3.25)]{Ba3}); \cite[(1.10)]{Ba3} with $m=2$, $E_v=2$ is the third graph of \cite[(3.25)]{Ba3}.
\end{proof}

\begin{lemma}[{$C$-parity; \cite[p.~434]{Ba3}: ``$\operatorname{tr}q^{2n+1}=0$''}]\label{sup4:lem:S212}
After S18--S19 every scalar line is a scalar kernel times $\mathbf 1_N$. If a graph without $R$-vertices has no external scalar leg and an odd number of external vector legs, its expression vanishes.
\end{lemma}

\begin{proof}
At \cite[(1.8)]{Ba3}/\cite[(1.10)]{Ba3} the two scalar legs are joined by $q^{m_v}$; at \cite[(1.6)]{Ba3} ($(\phi\cdot\phi)^2$) and \cite[(1.7)]{Ba3} the legs are paired by $\delta_{ab}$. With no external scalar legs the $\R^N$-index contractions close into loops, and each loop gives $\Tr\prod_vq^{m_v}=\Tr q^{M_\ell}$, $M_\ell$ the number of vector legs on the loop (all transporters are $1$ at $B_0=0$). $\sum_\ell M_\ell=\sum_vm_v=2I_v+E_v$ is odd, so some $M_\ell$ is odd, and $\Tr q^{\mathrm{odd}}=\Tr((q^{\mathrm{odd}})^T)=-\Tr q^{\mathrm{odd}}=0$.
\end{proof}

\begin{lemma}[{Taylor's formula \cite[(3.10)]{Ba3}}]\label{sup4:lem:S213}
For $f$ on $\eta\Z^d$ and a shortest (monotone) contour $\Gamma_{x,y}$ from $x$ to $y$,
\[
f(y)=f(x)+\sum_\mu(y_\mu-x_\mu)(\partial^\eta_\mu f)(x)+\sum_{b\subset\Gamma}s_b\eta\big[(\partial^\eta_{\mu_b}f)(b^{\mathrm{lo}})-(\partial^\eta_{\mu_b}f)(x)\big],
\]
$b^{\mathrm{lo}}$ the lower endpoint of $b$, $\mu_b$ its direction, $s_b=\operatorname{sgn}(y-x)_{\mu_b}$. The remainder is bounded by $|y-x|_1^{1+\alpha}$ times the $\alpha$-H\"older constant of $\partial f$ on the box spanned by $x$, $y$.
\end{lemma}

\begin{proof}
Telescoping along $\Gamma$: all $\mu$-bonds of a monotone $\Gamma$ have the same orientation, and there are $|y_\mu-x_\mu|/\eta$ of them. (Ba{\l}aban's $(b)_x$ is the forward bond at $x$ parallel to $b$.)
\end{proof}

\begin{lemma}[max-scale regrouping]\label{sup4:lem:S214}
Let $\cL$ be the line set of a graph, $I\subset\cL$ the internal lines of a subgraph $G_0$, and $\mathcal E\subset\cL\setminus I$, $\mathcal E\ne\emptyset$, its external lines (the lines of the graph with an end at a vertex of $G_0$). Fix a strict total order $\prec$ of $\cL$. For an index assignment $j\colon\cL\to\{0,\dots,k-1\}$ let $\tilde l(j)$ order $\cL$ by increasing $(j_l,l)$ lexicographically, and put $J(\tilde l)\defeq\{j:\tilde l(j)=\tilde l\}$; this is an admissible choice of the sets of \cite[(2.7)]{Ba3} (they partition all assignments, and $J(\tilde l)\subset\{j_{l(1)}\le\dots\le j_{l(m)}\}$). Let $\mathcal M$ be any set of orderings of $\cL\setminus I$ and $\mathcal L\defeq\{\tilde l:\tilde l|_{\cL\setminus I}\in\mathcal M$, and every $l\in I$ precedes every $e\in\mathcal E\}$. Then for every function $F$ of the assignment
\[
\sum_{\tilde l\in\mathcal L}\sum_{j\in J(\tilde l)}F(j)=\sum_{j_{\cL\setminus I}:\ \tilde l(j_{\cL\setminus I})\in\mathcal M}\ \sum_{j_I\in\prod_{l\in I}\{0,\dots,J_l-1\}}F(j),
\]
where $e^*\defeq\operatorname{argmin}_{e\in\mathcal E}(j_e,e)$ is the lexicographic minimum in $(j_e,\prec)$, which depends on $j_{\mathcal E}$, and $J_l\defeq j_{e^*}+[l\prec e^*]\in\{j_{e^*},j_{e^*}+1\}$. If $F$ is multilinear in the slice propagators of the lines of $I$, the inner sum replaces every $l\in I$ by the cumulative propagator $\sum_{i<J_l}G_{(i)}=\mathsf T_{J_l}^{-1}$ (\cref{sup4:cor:Z1}(i); for the photon lines the same with the vector covariance). This is the form of ``the same expression appears for all orderings \dots\ with the only condition that they are earlier than the external lines'' \cite[p.~437]{Ba3} and ``we sum over admissible orderings and indices'' \cite[p.~444]{Ba3}.
\end{lemma}

\begin{proof}
The $J(\tilde l)$ partition the assignments, so the left side is the sum over $\{j:\tilde l(j)\in\mathcal L\}$. The order induced on $\cL\setminus I$ depends on $j_{\cL\setminus I}$ only. ``Every $l\in I$ precedes every $e\in\mathcal E$'' is equivalent to $(j_l,l)<_{\mathrm{lex}}(j_{e^*},e^*)$ for all $l\in I$, and given $j_{\mathcal E}$ this is the product set $j_l\le J_l-1$.
\end{proof}

The lemma is a statement about sums over complete index ranges of the internal lines below the external scale; it asserts no identity for a fixed ordering or scale assignment.

\begin{lemma}[rescaling of the regrouped lines and coefficients]\label{sup4:lem:S215}
Let $j_0\defeq j_{e^*}$, $s\defeq L^{j_0}\eta$ and $\xi\defeq\eta/s=L^{-j_0}$. Rescaled by $s$ (\cref{sup4:cor:Z1}, Step~1), a regrouped scalar line $\mathsf T_{J_l}^{-1}$ becomes $s^{2-d}\breve G_l(\cdot/s,\cdot/s)$, where $\breve G_l$ is of the type of \cref{sup4:lem:R}: mass $m^2s^2\in[0,1]$, block term $a_{j_0}P$ on unit blocks if $J_l=j_0$, and $a_{j_0+1}L^{-2}P$ on $L$-blocks if $J_l=j_0+1$ (the case $J=k'+1$ of \cref{sup4:cor:Z1}). The same holds for the photon lines (componentwise; mass $\lmu^2t^2s^2$). A local coefficient of degree $D$ built from such lines (a sum over all but one vertex of a product of line kernels, derivatives and moments $(x'-x)^\beta$) equals $s^D$ times the same expression on the $\xi$-lattice; in $d=3$ with $D=0$ it is ``the same expression but on the scale $\xi$'' \cite[p.~444]{Ba3}.
\end{lemma}

\begin{proof}
Each line contributes $s^{2-d}$ times its rescaled kernel, each derivative $s^{-1}$, each moment factor $s$, each vertex sum $\eta^d=s^d\xi^d$ and each vertex $\eta$-power $s$; the total exponent is $D$ by \cite[(2.1)--(2.3)]{Ba3}.
\end{proof}

\begin{lemma}[the Ward--Takahashi identities, generated at full lines]\label{sup4:lem:S216}
Let $Y$ be $\xi\Z^d$ or a periodic box, $C\defeq C^\xi\otimes\mathbf 1_N$ (\cref{sup4:lem:C}), $C_{\mathrm{ph}}\defeq C^\xi\otimes\mathbf 1_d$, $A$ a finitely supported external bond field, and $F$ any polynomial in the $|\phi(x)|^2$ (for instance $\lambda\sum:\!|\phi|^4\!:$ and $\frac12\sum\delta m^2(x):\!|\phi(x)|^2\!:$, Wick-ordered with respect to any constants; mass counterterm vertices included). Put $Z(A)\defeq\int d\mu_{C_{\mathrm{ph}}}(A')\int d\phi\exp[-\frac12\ip\phi{(-\Delta^\xi_{A+A'}+1)\phi}-F(\phi)]$ and $W\defeq\log Z$ as formal power series in $(e,\lambda,\delta m^2,A)$. The Taylor coefficients of $W$ are the sums of all connected graphs of the given order and external structure with full lines $C$, $C_{\mathrm{ph}}$ (the vertices are those of \cite[(1.8)--(1.11)]{Ba3} with unrestricted $n$, $n'$, which agree with Ba{\l}aban's below order $\bar n+1$). Then: \textup{(a)} on a periodic box, $Z(A+\partial\lambda)=Z(A)$ for every periodic $\lambda$; \textup{(b)} every Taylor coefficient $W_n$ (symmetric $n$-linear in $A$ at fixed order in the couplings) obeys $W_n(\partial\lambda,A_2,\dots,A_n)=0$; \textup{(c)} (b) holds on $\xi\Z^d$ for finitely supported $\lambda$; \textup{(d)} the two-vector-leg kernel $\Pi_{\mu\nu}(y,y')$ of every order obeys $\sum_{y'}\xi^d\Pi_{\mu\nu}(y,y')=0$ for all $\mu,\nu,y$.
\end{lemma}

\begin{proof}
(a)--(c) are \cref{sup2:lem:S8cov,sup2:lem:S8torus,sup2:lem:S8lim} with $M^2=1$, the block term absent and $B=0$ (\cref{sup2:rem:S8scope}(c)): substituting $\phi=V_\lambda\psi$ has Jacobian $1$, $-\Delta_{A+\partial\lambda+A'}=V_\lambda(-\Delta_{A+A'})V_\lambda^*$ (\cref{sup4:lem:S171}(a); a periodic box has no boundary), and $F$ is invariant, which gives (a); comparing powers of $s$ in $W((sA)+s\partial\lambda)=W(sA)$ and polarising gives (b); at fixed order a coefficient is a finite sum of connected graphs whose lines are periodisations of the $\xi\Z^d$ lines (exponential decay, \cref{sup4:lem:C}), and with $A$, $\lambda$ finitely supported each graph sum converges absolutely to its $\xi\Z^d$ value as the box grows, which gives (c). (d) is \cref{sup2:lem:S8zero}(b) for $\Pi$: by (c), $\sum_{y',\nu}\xi^d\Pi_{\mu\nu}(y,y')(\lambda(y'+\xi e_\nu)-\lambda(y'))/\xi=0$ for all finitely supported $\lambda$, i.e.\ $\sum_\nu[\Pi_{\mu\nu}(y,z-\xi e_\nu)-\Pi_{\mu\nu}(y,z)]=0$ for every $z$; multiply by $z_\kappa$ and sum over $z$ (absolutely convergent by the exponential decay); the shift gives $\xi\sum_w\xi^d\Pi_{\mu\kappa}(y,w)=0$.
\end{proof}

\begin{lemma}[{lattice reflections; the symmetry behind \cite[(3.27)--(3.28), (3.37), (3.38)]{Ba3}}]\label{sup4:lem:S217}
For the theory of \cref{sup4:lem:S216} on $\xi\Z^d$ with constant coefficients, the reflection $R_\kappa$ ($x_\kappa\mapsto-x_\kappa$) acts by $\phi\mapsto\phi\circ R_\kappa$, $A_\mu(b)\mapsto A_\mu(R_\kappa b)$ ($\mu\ne\kappa$) and $A_\kappa(b)\mapsto-A_\kappa(b^r)$, where for $b=(x,x+\xi e_\kappa)$ we put $b^r\defeq(R_\kappa x-\xi e_\kappa,R_\kappa x)$. This leaves $Z$ invariant. Consequently, for the full fixed-order classes with lines $C$, $C_{\mathrm{ph}}$: \textup{(a)} the zero-momentum $\phi\phi A$ coefficient $\sum_{x_1,x_2}\xi^{2d}\Gamma_\kappa(x_1,x_2;b)$ vanishes; \textup{(b)} if $\sum_{y'}\xi^d\Pi_{\mu\nu}(y,y')=0$, then the first moment $\sum_{y'}\xi^d\Pi_{\mu\nu}(y,y')(y'_\lambda-y_\lambda)$ vanishes for all $\mu,\nu,\lambda$ ($d\ge2$).
\end{lemma}

\begin{proof}
\emph{Invariance.} $|U(A_b)\phi(b_+)-\phi(b_-)|=|U(-A_b)\phi(b_-)-\phi(b_+)|$ by orthogonality, so reversing a bond with $A\mapsto-A$ preserves the action; $C$ is even, and $C_{\mathrm{ph}}$ is invariant under $A_\kappa\mapsto-A_\kappa$. \emph{(a)} $\Gamma_\kappa(R_\kappa x_1,R_\kappa x_2;b^r)=-\Gamma_\kappa(x_1,x_2;b)$, and by translation invariance $\sum_{x_1,x_2}\Gamma_\kappa(x_1,x_2;b)$ does not depend on $b$; so the sum equals minus itself. \emph{(b)} Write $m(b)$ for the bond midpoint. Translation invariance gives $\Pi_{\mu\nu}(b,b')=\pi_{\mu\nu}(m(b')-m(b))$ and the invariance gives $\pi_{\mu\nu}(R_\kappa u)=s^\kappa_\mu s^\kappa_\nu\pi_{\mu\nu}(u)$ with $s^\kappa_\mu=-1$ iff $\mu=\kappa$. Since $y'_\lambda-y_\lambda=(m(b')-m(b))_\lambda+\frac\xi2(\delta_{\mu\lambda}-\delta_{\nu\lambda})$ and $\sum_u\pi_{\mu\nu}(u)=0$ by hypothesis, the first moment is $M\defeq\sum_u\xi^d\pi_{\mu\nu}(u)u_\lambda$. Reflection in $\kappa=\lambda$ gives $M=-s^\lambda_\mu s^\lambda_\nu M$, so $M=0$ unless exactly one of $\mu,\nu$ equals $\lambda$; reflection in $\kappa\ne\lambda$ gives $M=s^\kappa_\mu s^\kappa_\nu M$, so $M=0$ unless, for every $\kappa\ne\lambda$, $\mu=\kappa\Leftrightarrow\nu=\kappa$. Both conditions together are impossible: if $\mu=\lambda\ne\nu$, then $\kappa=\nu$ violates the second.
\end{proof}

\cite[(3.27)--(3.28)]{Ba3} (a computation for the one-loop bubble) and \cite[(3.37), (3.38)]{Ba3} (``the functions which we are summing are odd'') are instances of (b) and (a).
